% =============================================================================
% Wprowadzenie. Plik edytowany ręcznie -- skrypt analiza_odpowiedzi_naukowa.py go NIE nadpisuje.
% Liczby: definicje.tex; tabele: \TabDir (CSV); rysunki: \FigDir; literatura: references.bib.
% =============================================================================
\section{Wprowadzenie}

Przypominanie informacji z pamięci utrwala je silniej niż ponowne czytanie (efekt testowania),
a nauka rozłożona w czasie jest trwalsza niż nauka skumulowana (efekt rozłożenia w czasie)
\citep{roediger2006,cepeda2006}. Oba efekty łączy edukacja rozłożona w czasie (ang. \emph{spaced
education}), w której krótkie pytania wielokrotnego wyboru wracają cyklicznie przez tygodnie lub
miesiące. W badaniach z udziałem lekarzy i innych pracowników ochrony zdrowia poprawiała ona
przyswajanie i utrwalanie wiedzy \citep{kerfoot2009,phillips2019}.

W organizacjach podobną formę ma \emph{mikrouczenie} (ang. \emph{microlearning}), czyli krótkie
jednostki szkoleniowe, często w formie quizów. Dowody jego skuteczności pochodzą jednak głównie
z badań na małych próbach i z tzw. szarej literatury \citep{taylor2022}, a efekty obowiązkowych
szkoleń pracowniczych, np.\ z bezpieczeństwa informacji, bywają niewielkie lub zanikają z czasem
\citep{reinheimer2020}. Skuteczność szkoleń ocenia się zwykle przez transfer, czyli zastosowanie
wiedzy w pracy \citep{blume2010}. Transfer wymaga jednak, aby wiedza została najpierw nabyta
i utrwalona.

Ten wcześniejszy etap można sprawdzić na danych z platform mikrouczenia. Platforma zapisuje każdą
odpowiedź wraz z czasem i wynikiem, a to samo pytanie wraca do tego samego pracownika wielokrotnie.
Kolejne odpowiedzi stanowią więc powtarzane pomiary tej samej wiedzy u tej samej osoby. Podobne dane
wykorzystuje się w modelach ucznia, stosowanych w systemach tutorskich i aplikacjach do nauki języków
\citep{corbett1994,settles2016}.

Badanie przeprowadzono w banku, w którym pracownicy odpowiadali na krótkie pytania wyświetlane na
komputerze w trakcie dnia pracy, a system gromadził odpowiedzi przez trzy lata. Pytania wracały według
harmonogramu niezależnego od wyniku odpowiedzi, a po każdej odpowiedzi wyświetlano odpowiedź poprawną.
Taki harmonogram pozwala porównać korekty błędów i zapominanie bez selekcji, którą wprowadzają
algorytmy adaptacyjne, co w danych z platform szkoleniowych występuje rzadko. Badania laboratoryjne
wskazują ponadto, że wybrany dystraktor może utrwalić się w pamięci, a błędy popełnione z wysoką
pewnością mogą powracać mimo informacji zwrotnej \citep{butler2008,butler2011}. Nie wiadomo, czy
zjawisko to występuje w szkoleniach pracowniczych. Sformułowano trzy pytania badawcze:
(PB1) czy odsetek poprawnych odpowiedzi rośnie, gdy to samo pytanie wraca do tego samego pracownika;
(PB2) jak poprawność zależy od wyniku poprzedniej próby i od czasu, który od niej upłynął;
(PB3) czy powtórzone błędy utrzymują się mimo informacji zwrotnej, co wskazywałoby raczej na błędne
przekonania niż na zgadywanie.
